% problem_id: S001-lemma-galois-inertia
% source_id: S001 (Stacks Project)
% source_file: more-algebra.tex
% statement_locator: more-algebra.tex lines 33250--33281
% proof_locator: more-algebra.tex lines 33283--33406
% env: lemma
% id_note: label 在本来源内唯一
% pairing: tight (verified)
% status: complete (statement + author proof verbatim + reason + locators)
%
\begin{lemma}
\label{lemma-galois-inertia}
Let $A$ be a discrete valuation ring with fraction field $K$.
Let $L/K$ be a finite Galois extension with Galois group $G$.
Let $B$ be the integral closure of $A$ in $L$. Let $\mathfrak m \subset B$
be a maximal ideal. The inertia group $I$ of $\mathfrak m$
sits in a canonical exact sequence
$$
1 \to P \to I \to I_t \to 1
$$
such that
\begin{enumerate}
\item if $D$ is the decomposition group we have
$$
P = \{\sigma \in I \mid \sigma|_{\mathfrak m/\mathfrak m^2} =
\text{id}_{\mathfrak m/\mathfrak m^2}\}
=
\{\sigma \in D \mid \sigma \text{ acts trivially on }\text{Gr}_\mathfrak m(B)\}
$$
\item $P$ is a normal subgroup of $D$,
\item $P$ is a $p$-group if the characteristic of $\kappa_A$ is
$p > 0$ and $P = \{1\}$ if the characteristic of $\kappa_A$ is zero,
\item $I_t$ is cyclic of order the prime to $p$ part of the integer $e$,
\item there is a canonical isomorphism
$\theta : I_t \to \mu_e(\kappa(\mathfrak m))$, and
\item
$P =
\{\sigma \in D \mid \sigma|_{B/\mathfrak m^2} = \text{id}_{B/\mathfrak m^2}\}$
if $\kappa(m)$ is separable over the residue field of $A$.
\end{enumerate}
Here $e$ is the integer of Lemma \ref{lemma-galois-conclusion}.
\end{lemma}

\begin{proof}
Recall that $|G| = [L : K] = nef$, see Lemma \ref{lemma-galois-conclusion}.
Since $G$ acts transitively on the set
$\{\mathfrak m_1, \ldots, \mathfrak m_n\}$ of maximal ideals of $B$
(Lemma \ref{lemma-galois})
and since $D$ is the stabilizer of an element we see that $|D| = ef$.
By Lemma \ref{lemma-galois-galois} we have
$$
ef = |D| = |I| \cdot |\text{Aut}(\kappa(\mathfrak m)/\kappa)|
$$
where $\kappa$ is the residue field of $A$.
As $\kappa(\mathfrak m)$ is normal over $\kappa$ the order of
$\text{Aut}(\kappa(\mathfrak m)/\kappa)$ differs from $f$ by
a power of $p$ (see
Fields, Lemma \ref{fields-lemma-normal-and-automorphisms}
and discussion following
Fields, Definition \ref{fields-definition-insep-degree}).
Hence the prime to $p$ part
of $|I|$ is equal to the prime to $p$ part of $e$.

\medskip\noindent
Set $C = B_\mathfrak m$. Then $I$ acts on $C$ over $A$ and trivially
on the residue field of $C$. Let $\pi_A \in A$ and $\pi_C \in C$ be
uniformizers. Write $\pi_A = u \pi_C^e$ for some unit $u$ in $C$.
For $\sigma \in I$ write $\sigma(\pi_C) = \theta_\sigma \pi_C$ for some
unit $\theta_\sigma$ in $C$. Then we have
$$
\pi_A = \sigma(\pi_A) = \sigma(u) (\theta_\sigma \pi_C)^e
= \sigma(u) \theta_\sigma^e \pi_C^e = \frac{\sigma(u)}{u} \theta_\sigma^e \pi_A
$$
Since $\sigma(u) \equiv u \bmod \mathfrak m_C$ as $\sigma \in I$
we see that the image $\overline{\theta}_\sigma$
of $\theta_\sigma$ in $\kappa_C = \kappa(\mathfrak m)$
is an $e$th root of unity.
We obtain a map
\begin{equation}
\label{equation-inertia-character}
\theta : I \longrightarrow \mu_e(\kappa(\mathfrak m)),\quad
\sigma \mapsto \overline{\theta}_\sigma
\end{equation}
We claim that $\theta$ is a homomorphism of groups and independent
of the choice of uniformizer $\pi_C$. Namely, if $\tau$ is a second
element of $I$, then
$\tau(\sigma(\pi_C)) = \tau(\theta_\sigma \pi_C) =
\tau(\theta_\sigma) \theta_\tau \pi_C$, hence
$\theta_{\tau \sigma} = \tau(\theta_\sigma) \theta_\tau$ and
since $\tau \in I$ we conclude that
$\overline{\theta}_{\tau \sigma} =
\overline{\theta}_\sigma \overline{\theta}_\tau$.
If $\pi'_C$ is a second uniformizer, then we see
that $\pi'_C = w \pi_C$ for some unit $w$ of $C$ and
$\sigma(\pi'_C) = w^{-1}\sigma(w)\theta_\sigma \pi'_C$,
hence $\theta'_\sigma = w^{-1}\sigma(w)\theta_\sigma$,
hence $\theta'_\sigma$ and $\theta_\sigma$
map to the same element of the residue field as before.

\medskip\noindent
Since $\kappa(\mathfrak m)$ has characteristic $p$, the group
$\mu_e(\kappa(\mathfrak m))$ is cyclic of order at most the prime
to $p$ part of $e$ (see Fields, Section \ref{fields-section-roots-of-1}).

\medskip\noindent
Let $P = \Ker(\theta)$. By construction the elements of $P$ are exactly
the elements of $I$ which act trivially on
$\mathfrak m/\mathfrak m^2 = \text{Gr}_\mathfrak m^1(B)$,
i.e., $P = \{\sigma \in I \mid \sigma|_{\mathfrak m/\mathfrak m^2} =
\text{id}_{\mathfrak m/\mathfrak m^2}\}$. Also
$I$ consists of the elements of $D$ which act trivially on
$\kappa(\mathfrak m) = B/\mathfrak m$.
Since the graded ring $\text{Gr}_\mathfrak m(B)$ is generated
by $\text{Gr}^1_\mathfrak m(B) = \mathfrak m/\mathfrak m^2$ over
$\text{Gr}^0_\mathfrak m(B) = B/\mathfrak m$, we conclude that
$P = \{\sigma \in D \mid \sigma \text{ acts trivially on }
\text{Gr}_\mathfrak m(B)\}$.
Thus (1) is true. This implies (2) as $P$ is the kernel
of the homomorphism $D \to \text{Aut}(\text{Gr}_\mathfrak m(B))$.
If we can prove (3), then parts (4) and (5) will follow as $I_t$
will be isomorphic to $\mu_e(\kappa(\mathfrak m))$ as the arguments above show
that $|I_t| \geq |\mu_e(\kappa(\mathfrak m))|$.

\medskip\noindent
Thus it suffices to prove that the
kernel $P$ of $\theta$ is a $p$-group.
Let $\sigma$ be a nontrivial element of the kernel. Then $\sigma - \text{id}$
sends $\mathfrak m_C^i$ into $\mathfrak m_C^{i + 1}$
for all $i$. Let $m$ be the order of $\sigma$. Pick $c \in C$ such
that $\sigma(c) \not = c$. Then $\sigma(c) - c \in \mathfrak m_C^i$,
$\sigma(c) - c \not \in \mathfrak m_C^{i + 1}$ for some $i$ and
we have
\begin{align*}
0
& =
\sigma^m(c) - c \\
& =
\sigma^m(c) - \sigma^{m - 1}(c) + \ldots + \sigma(c) - c \\
& =
\sum\nolimits_{j = 0, \ldots, m - 1} \sigma^j(\sigma(c) - c) \\
& \equiv
m(\sigma(c) - c) \bmod \mathfrak m_C^{i + 1}
\end{align*}
It follows that $p | m$ (or $m = 0$ if $p = 1$). Thus every element of the
kernel of $\theta$ has order divisible by $p$, i.e., $\Ker(\theta)$
is a $p$-group.

\medskip\noindent
Proof of (6). Assume $\kappa(\mathfrak m)/\kappa$ is separable.
In this case $f = |D/I|$ by Lemma \ref{lemma-galois-galois}
and $I$ has order $e$. If $e = 1$, then $P = I = \{\text{id}\}$
and the result is true. If $e > 1$, then
$B/\mathfrak m^2 = C/\pi_C^2C$ is a $\kappa$-algebra
and a small extension of $\kappa(\mathfrak m)$.
Because $\kappa(\mathfrak m)$ is formally \'etale over $\kappa$
(Algebra, Lemma
\ref{algebra-lemma-characterize-separable-algebraic-field-extensions})
there is a unique $\kappa$-algebra map
$\kappa(\mathfrak m) \to B/\mathfrak m^2$
right inverse to $B/\mathfrak m^2 \to \kappa(\mathfrak m)$.
In other words, there is a $D$-equivariant isomorphism
$B/\mathfrak m^2 = \mathfrak m/\mathfrak m^2 \oplus \kappa(\mathfrak m)$
of $\kappa$-algebras. This immediately shows that
$P =
\{\sigma \in D \mid \sigma|_{B/\mathfrak m^2} = \text{id}_{B/\mathfrak m^2}\}$
in this case.
\end{proof}
